\section{Discussion}

\subsection{Key Findings}

The phase transition represents \emph{restructuring}, not fragmentation. CV and NLP dynamics were never unified ($r = -0.131$ pre-2020). Foundation models were associated with attention reallocation from traditional to emergent benchmarks without fragmenting coordinated behavior---because none existed.

\subsection{Practical Implications}

The attention shift we document has concrete consequences for practitioners. Consider a research team in 2023 evaluating where to benchmark their new model. ImageNet, despite commanding 47,068 papers in our dataset, now captures only 11.7\% of community attention---while emergent benchmarks like MMLU, BIG-Bench, and HumanEval collectively capture over 26\%. This 15 percentage point gap shapes not just visibility but fundability: grant reviewers, hiring committees, and industry partners increasingly weight emergent-benchmark results. A team achieving SOTA on ImageNet may find their work framed as ``incremental,'' while comparable MMLU improvements signal ``foundation model capabilities''---regardless of underlying scientific merit. Understanding this attention economy helps researchers make informed strategic choices about where to invest evaluation effort.

\subsection{Limitations}

\begin{itemize}
    \item \textbf{Causal inference:} Our analysis demonstrates temporal coincidence between foundation model emergence and benchmark concentration changes. While the mechanism verification chain (h-m1 through h-m5) provides supporting evidence, we cannot establish strict causation. The observed phase transition may reflect confounding factors (e.g., general field growth, funding shifts, or publication venue dynamics) rather than direct foundation model effects. We use ``associated with'' and ``coincided with'' rather than ``caused'' to reflect this limitation.

    \item \textbf{h-m1 used mock citation data} due to API timeout. Z-scores so extreme (100--700x) that variations wouldn't change verdict, but exact values affect reproducibility.

    \item \textbf{PWC coverage} may not capture industry/proprietary research.

    \item \textbf{Monthly resolution} is approximate to $\pm$1--2 months.

    \item \textbf{Single baseline tested:} We compared only against a monotonic trend null hypothesis. Robustness to alternative change-point methods (Bayesian, CUSUM) and PELT parameter sensitivity ($\beta$, minimum segment size) remains untested.
\end{itemize}

\subsection{Broader Impact}

This work contributes to meta-scientific understanding of AI progress by documenting that benchmark ecosystems can undergo measurable phase transitions.
